\section{Základní části počítače}\label{sec:hw-zakladni-casti}

Slovem \emph{hardware} označujeme fyzické součásti počítače: procesor, paměťové čipy, základní desku, úložiště, klávesnici, displej a~další zařízení. Samotný hardware však neurčuje, jakou konkrétní úlohu má počítač řešit. Potřebný postup popisuje \emph{software}, tedy programy a~data, s~nimiž tyto programy pracují. Stejný počítač proto může chvíli zobrazovat webovou stránku, poté překládat program a~nakonec přehrávat video. Hardware zůstává stejný, mění se pouze instrukce, které právě vykonává.

Při velmi hrubém pohledu můžeme činnost počítače rozdělit mezi následující části:
\begin{itemize}
    \item \textbf{datová část} provádí výpočty a~zpracovává data;
    \item \textbf{řídicí část} určuje, co a~v~jakém pořadí se má stát;
    \item \textbf{paměť} uchovává programy, jejich instrukce, vstupní data a~mezivýsledky;
    \item \textbf{vstupní zařízení} předávají počítači data z~okolí;
    \item \textbf{výstupní zařízení} předávají výsledky uživateli nebo jinému zařízení.
\end{itemize}

Datová a~řídicí část jsou v~běžném počítači součástí \emph{procesoru}, často označovaného zkratkou CPU z~anglického \emph{Central Processing Unit}. Procesor je sice místem, kde se instrukce skutečně provádějí, bez paměti a~vstupně/výstupních zařízení by však neměl odkud instrukce ani data získat a~kam uložit výsledky. Jednotlivé části proto nelze chápat jako izolované krabičky; při vykonávání programu spolu neustále komunikují.

\begin{figure}[ht]
    \centering
    \begin{tikzpicture}[
    x=1cm,
    y=1cm,
    >=Stealth,
    every node/.style={font=\small},
    block/.style={
        draw,
        thick,
        rounded corners,
        align=center,
        minimum height=0.95cm
    },
    inner/.style={
        draw,
        rounded corners,
        fill=white,
        align=center,
        minimum width=3.1cm,
        minimum height=0.55cm
    },
    arrow/.style={->, thick},
    both/.style={<->, thick}
]
    \node[
        block,
        fill=blue!10,
        minimum width=4.4cm,
        minimum height=4.1cm
    ] (cpu) at (0,0.5) {};
    \node[font=\large\bfseries] at (0,1.95) {Procesor (CPU)};
    \node[inner] at (0,1.15) {aritmeticko-logická jednotka};
    \node[inner] at (0,0.35) {řadič};
    \node[inner] at (0,-0.45) {registry};

    \node[
        block,
        fill=green!12,
        minimum width=3.5cm
    ] (memory) at (0,-3.0) {\textbf{Paměť}\\[-1mm]\scriptsize instrukce a~data};

    \node[
        block,
        fill=orange!15,
        minimum width=2.7cm,
        minimum height=1.15cm
    ] (input) at (-5.3,0.5) {\textbf{Vstupní}\\ zařízení};

    \node[
        block,
        fill=orange!15,
        minimum width=2.7cm,
        minimum height=1.15cm
    ] (output) at (5.3,0.5) {\textbf{Výstupní}\\ zařízení};

    \draw[arrow] (input.east) -- node[above] {data} (cpu.west);
    \draw[arrow] (cpu.east) -- node[above] {výsledky} (output.west);
    \draw[both] (cpu.south) -- node[right, align=left] {data\\ instrukce} (memory.north);
\end{tikzpicture}

    \caption{Zjednodušené schéma základních částí počítače.}
    \label{fig:hw-zakladni-casti}
\end{figure}

Na obrázku~\ref{fig:hw-zakladni-casti} je zachycen typický tok informací. Vstupní zařízení dodá data, procesor je podle instrukcí programu zpracuje a~výsledek může předat výstupnímu zařízení. Během tohoto procesu se instrukce i~data přesouvají mezi procesorem a~pamětí. Rozdělení na vstup a~výstup přitom není vždy absolutní. Například síťová karta data přijímá i~odesílá a~dotyková obrazovka současně zobrazuje výstup a~zaznamenává vstup uživatele.

\notebox{V~této kapitole se zaměříme hlavně na procesor a~paměťový subsystém, tedy na části, které přímo souvisejí s~vykonáváním programu. Vstupně/výstupní zařízení budeme uvažovat pouze na úrovni potřebné k~pochopení celkového schématu.}

\paragraph{Data, instrukce a~jejich význam.}

Uvnitř počítače jsou informace reprezentovány pomocí bitů. Samotná posloupnost bitů však ještě neříká, co představuje. Stejný vzor může být podle okolností interpretován jako celé číslo, část obrázku, znak textu, adresa v~paměti nebo strojová instrukce. Význam tedy nevychází pouze z~uložených bitů, ale také z~toho, jak s~nimi právě pracuje procesor a~jaký formát programu očekává.

Tento princip je důležitý i~pro programování. Proměnná určitého datového typu určuje, jak má program uloženou hodnotu chápat, zatímco instrukce určuje, jaká operace se má s~hodnotou provést. Na úrovni hardware jsou však obě informace stále reprezentovány bity a~pro jejich přenos se používají stejné základní fyzikální prostředky.

\tipbox{Při hledání chyby v~programu je užitečné rozlišovat mezi uloženou hodnotou a~její interpretací. Bity mohou být přeneseny nebo uloženy správně, ale program je může vyložit v~nesprávném formátu.}
